\documentclass[conference]{IEEEtran}
\usepackage[utf8]{inputenc}
\usepackage{amsmath,amssymb}
\usepackage{microtype}
\usepackage{booktabs}
\usepackage{graphicx}
\usepackage{xcolor}
\usepackage{array}
\usepackage{hyperref}
\hypersetup{colorlinks=true, urlcolor=blue, linkcolor=black, citecolor=black}
\usepackage{cite}

% Compact float, list, and line spacing to guarantee strict 4-page IEEE fit in Overleaf
\renewcommand{\baselinestretch}{0.925}
\setlength{\textfloatsep}{2.5pt plus 1.0pt minus 1.0pt}
\setlength{\floatsep}{3.5pt plus 1.0pt minus 1.0pt}
\setlength{\dbltextfloatsep}{2.5pt plus 1.0pt minus 1.0pt}
\setlength{\dblfloatsep}{3.5pt plus 1.0pt minus 1.0pt}
\setlength{\abovecaptionskip}{2.0pt}
\setlength{\belowcaptionskip}{0pt}
\AtBeginDocument{%
  \abovedisplayskip=4.0pt plus 1.0pt minus 2.0pt%
  \belowdisplayskip=4.0pt plus 1.0pt minus 2.0pt%
}

% Helper macro to load project figures via relative path or local directory fallback
\newcommand{\incplot}[3]{%
  \IfFileExists{#1}{\includegraphics[width=#3]{#1}}{%
    \IfFileExists{#2}{\includegraphics[width=#3]{#2}}{%
      \fbox{\parbox[c][1.1in][c]{#3}{\centering\scriptsize\texttt{\detokenize{#2}}}}%
    }%
  }%
}

% Abstract styling for IEEE conference format
\makeatletter
\def\abstract{\normalfont
  \if@twocolumn
    \@IEEEabskeysecsize\bfseries\textit{\abstractname}---\relax
  \else
    \bgroup\par\addvspace{0.5\baselineskip}\centering\vspace{-1.78ex}\@IEEEabskeysecsize\textbf{\abstractname}\par\addvspace{0.5\baselineskip}\egroup\quotation\@IEEEabskeysecsize
  \fi\@IEEEgobbleleadPARNLSP}
\makeatother

\title{\vspace{-12pt}A Queueing-Network Simulator for an LLM Inference Serving Pipeline: Phase-2 Progress Report}

\author{
  \IEEEauthorblockN{Vaibhav Mahore (23631)\IEEEauthorrefmark{1}, Sujal Chavda (24288)\IEEEauthorrefmark{1}, Arindam Raina (25880)\IEEEauthorrefmark{1}, Shreemaye Das (25958)\IEEEauthorrefmark{2}}
  \IEEEauthorblockA{\IEEEauthorrefmark{1}Department of Computer Science and Automation (CSA), Indian Institute of Science, Bangalore, India}
  \IEEEauthorblockA{\IEEEauthorrefmark{2}Department of Computational and Data Sciences (CDS), Indian Institute of Science, Bangalore, India}
  \IEEEauthorblockA{Course: E0 240 (Modeling and Simulation) \quad Email: \{mvaibhav, sujals, arindamr, shreemayed\}@iisc.ac.in}
}

\begin{document}
\maketitle\vspace{-14pt}

\begin{abstract}
Large language model (LLM) inference serving systems must satisfy strict Time-to-First-Token (TTFT) and Inter-Token Latency (ITL) Service Level Objectives (SLOs) under bursty arrivals and variable generation lengths. We design and implement a cycle-driven discrete-time queueing-network simulator that models the serving pipeline as an open three-station tandem network: a deterministic admission and tokenization server ($Q_1$: $Geo/D/1$, $S_1{=}2$), a dual-trigger First-Come First-Served (FCFS) batching scheduler ($Q_2$: cap $B$, timeout $\tau$), and a parallel multi-worker GPU decode pool ($Q_3$: $Geo/Geo/c$ or sublinear batch service $S_3(b)$). Directly addressing Phase-1 feedback, we establish a structural systems-to-queueing mapping that decomposes TTFT and ITL into station-level waiting delays ($W_{q,1}, W_{q,2}, W_{q,3}$) and service times. We then execute a two-layer validation and experimental study. In Layer~1, we verify our engine against discrete-time closed forms over $10^7$ cycles, achieving worst-case occupancy errors of $0.75\%$ ($Geo/D/1$, $S_1{=}2$), $0.46\%$ ($S_1{=}3$), and $0.94\%$ ($Geo/Geo/1$, $E[S]{=}2$), reproducing the Pollaczek-Khinchine factor-2 variance penalty within $0.31\%$ at $\lambda{=}0.45$ and network flow conservation within $3.4{\times}10^{-7}$. In Layer~2, we ground service scaling in batched inference sweeps across $B \in \{1, 2, 4, 8, 16, 32\}$ on an Nvidia Tesla T4 GPU across four model families (\texttt{Qwen2.5-1.5B}, \texttt{Gemma-2-2B}, \texttt{Llama-3.2-1B}, and \texttt{Phi-3-mini-4k}). Calibrating a power-law model on \texttt{Qwen2.5-1.5B} ($B \in \{1, 4, 16\}$) predicts held-out sizes $\{2, 8, 32\}$ with $4.43\%$ latency MAPE and $4.18\%$ throughput MAPE, while an uncapped 100-prompt diagnostic (512-token ceiling) reveals non-geometric output lengths ($KS{=}0.3965$) where straggler length $\max_{j \le B} L_j$ governs batch runtime ($r{=}0.9493$, cutting uncapped latency MAPE from $15.25\%$ to $9.14\%$). Finally, experiments E2 and E3 prove that $c{=}2$ parallel workers absorb an $80.2\%$ service-variability penalty at $\lambda{=}0.45$, while static batching ($\tau{=}8, B{\ge}8$) cuts GPU worker utilization by $48.9\%$ ($0.5999 \to 0.3063$).
\end{abstract}

\begin{figure*}[!t]
\centering
\incplot{fig_arch_mapping.pdf}{figures/fig_arch_mapping.pdf}{0.76\textwidth}
\vspace{-4pt}
\caption{Structural mapping from a physical LLM serving stack (amber strips) to our three-station discrete-time queueing network (blue bodies), addressing Phase-1 Feedback~\#1. Squares are requests, rounded boxes are batches and circles are servers. The bar splits a request's delay into queueing waits $W_{q,i}$ and service times (not to scale) and shows how TTFT, ITL and end-to-end sojourn ($W_{\text{soj}}$) are composed.}
\label{fig:arch_mapping}
\vspace{-6pt}
\end{figure*}

\section{Problem Statement and Motivation}
Serving generative Large Language Models (LLMs) such as LLaMA and Qwen in interactive applications requires satisfying strict latency Service Level Objectives (SLOs) under stochastic request arrivals~\cite{kwon2023vllm,agrawal2024vidur}: \emph{Time-to-First-Token} (TTFT), measuring initial responsiveness, and \emph{Inter-Token Latency} (ITL), measuring token generation speed.

\subsection{Why Model LLM Serving as a Queueing Network? (Addressing Phase-1 Feedback \#1)}
A central question raised in our Phase-1 evaluation is: \emph{why should an LLM inference pipeline be modeled as a network of queues rather than analyzed as a single compute kernel or evaluated solely through full-system profiling?} We establish three foundational justifications, illustrated in Fig.~\ref{fig:arch_mapping}:

\textbf{1) Physical Decoupling of Asynchronous Resource Stages.}
In production serving stacks (vLLM~\cite{kwon2023vllm}, Hugging Face TGI, TensorRT-LLM), requests traverse three decoupled stages governed by distinct queueing mechanics:
\begin{itemize}\setlength{\itemsep}{0pt}\setlength{\topsep}{1pt}
  \item \emph{CPU Admission \& Tokenization ($Q_1$):} An asynchronous front-end parses HTTP/gRPC payloads and executes Byte-Pair Encoding (BPE) tokenization on the CPU. Because tokenization takes a small, deterministic slice ($0.206\text{ ms}$ mean in our T4 host trials) decoupled from GPU contention, it forms a single-server deterministic queue ($Geo/D/1$).
  \item \emph{Scheduler Batch Accumulation ($Q_2$):} Autoregressive decoding is memory-bandwidth bound on GPU High-Bandwidth Memory (HBM); serving $B{=}1$ wastes compute capacity. The scheduler therefore buffers tokenized requests in an FCFS queue until $B$ requests accumulate or a latency timer $\tau$ expires~\cite{yu2022orca}, introducing an explicit \emph{batch-aggregation delay} $W_{q,2}$.
  \item \emph{Parallel GPU Execution Contention ($Q_3$):} Formed batches are dispatched to $c$ parallel GPU execution streams. Because output token lengths vary widely across prompts ($83\text{--}487$ tokens in our uncapped diagnostic), active batches occupy workers for stochastic durations, subjecting arriving batches to \emph{worker contention delay} $W_{q,3}$.
\end{itemize}

\textbf{2) Exact Additive Decomposition of User-Perceived Latency.}
End-to-end request sojourn time $W_{\text{soj}}$ decomposes additively across all three stations:
\begin{equation}
W_{\text{soj}} = \underbrace{W_{q,1} + W_{q,2} + W_{q,3}}_{W_{\text{e2e,qwait}}\text{ (Queueing Delay)}} + \underbrace{S_1 + S_3(b, L_{\max})}_{\text{Service Duration}}.
\label{eq:sojourn_decomp}
\end{equation}
As load approaches capacity ($\rho_i \to 1$), queueing delay $W_{\text{e2e,qwait}}$ dominates hardware compute time by over $15\times$ (growing from $0.057$ cycles at $\lambda{=}0.05$ to $65.15$ cycles at $\lambda{=}0.49$ in Experiment~E2). A multi-station network isolates whether tail latency stems from admission bottlenecks ($Q_1$), overly patient batching timers ($\tau$ at $Q_2$), or decode-length variance amplification ($Q_3$).

\textbf{3) Analytical Tractability and Speed vs.\ Operator Profiling.}
Trace-driven simulators like Vidur~\cite{agrawal2024vidur} profile thousands of tensor operators across GPU clusters, requiring costly hardware access and offering no closed-form insight into stability limits. Our discrete-time queueing simulator evaluates $10^7$ cycles in $\approx 6.5\text{ s}$ on a single CPU core, sweeping $(\lambda, B, \tau, c)$ analytically and numerically while grounding service laws in real Tesla T4 GPU measurements.

\section{Related Work}
\textbf{LLM Serving Systems.} Orca~\cite{yu2022orca} pioneered iteration-level (continuous) scheduling to retire completed sequences at token boundaries. vLLM~\cite{kwon2023vllm} introduced PagedAttention to eliminate KV-cache fragmentation, unlocking larger batch caps $B$. FlexGen~\cite{sheng2023flexgen} maximized single-GPU offline throughput via CPU/disk offloading. While these systems engineer runtime mechanisms, our work formulates a verifiable queueing model of their batching and worker-pool trade-offs.

\textbf{Simulators and Queueing Theory.} Vidur~\cite{agrawal2024vidur} predicts cluster scheduling latency via machine-learned operator tables but lacks closed-form queueing verification. Following E0~240's discrete-time modeling principles~\cite{mor2006discrete}, we verify our baseline nodes against exact $Geo/D/1$ and $Geo/Geo/1$ closed forms and calibrate $Q_3$ using empirical Tesla T4 traces~\cite{zheng2023response}.

\section{Methodology and Simulator Architecture}

\subsection{Three-Station Tandem Queueing Formulation}
Time advances in discrete cycles $t \in \{1, \dots, T\}$. External arrivals follow a Bernoulli process with rate $\lambda \in (0, 0.5)$ requests/cycle.
\textbf{Station 1 ($Q_1$):} Single-server FIFO $Geo/D/1$ queue with deterministic service $S_1 = 2$ cycles ($\mu_1 = 0.5$; $S_1 = 3$ also verified), using the course's same-cycle immediate-access convention (\texttt{single\_queue\_simulation\_geo\_d\_1.m}).
\textbf{Station 2 ($Q_2$):} FCFS batch buffer holding $q_2(t)$ requests with head-of-line arrival cycle $t_{\text{head}}$. At cycle $t$, $Q_2$ releases a batch of size $b$ to $Q_3$ when either trigger fires:
\begin{equation}
q_2(t) \ge B \implies b = B, \quad \text{or} \quad t - t_{\text{head}} \ge \tau \implies b = q_2(t) < B.
\end{equation}
With $B{=}1$ or $\tau{=}0$, $Q_2$ is a zero-delay pass-through ($W_{q,2}{=}0$).
\textbf{Station 3 ($Q_3$):} Pool of $c$ homogeneous parallel workers, each serving one entire batch synchronously (all requests in the batch share the batch service duration) under either: (i)~a \emph{deterministic sublinear law} $S_3(b) = \lceil S_{\text{base}} + \delta(b - 1) \rceil$ ($S_{\text{base}}{=}2$), or (ii)~a \emph{geometric law} with completion probability $p_s(b) = 1/(S_{\text{base}} + \delta(b - 1))$ and mean $E[S_3(b)] = 1/p_s(b)$.

\textbf{Stability Criteria.} Because requests are conserved along the path ($\lambda_1 = \lambda_2 = \lambda_3 = \lambda$), the binding stability conditions are:
\begin{equation}
\rho_1 = \lambda S_1 = 2\lambda < 1 \implies \lambda < 0.50, \qquad \rho_3 = \frac{\lambda E[S_3]}{c} < 1,
\end{equation}
where batch service duration $S_3(b)$ depends on the realized batch size $b$, and $Q_2$ is stable whenever $Q_3$ is stable.

\subsection{Cycle-Driven Engine and Common Random Numbers (CRN)}
In each cycle $t$, \texttt{network\_sim.py} executes four steps in exact class-code order: (1)~decrement deterministic countdowns or draw geometric completions via Inverse Transform sampling ($U \le p_s$); (2)~route completed jobs downstream ($Q_1 \to Q_2 \to Q_3 \to \text{Exit}$); (3)~check $Q_2$ release triggers and dispatch queued jobs to idle workers; and (4)~accumulate post-dispatch queue lengths $Q_i(t)$ (excluding in-service requests) and busy worker counts $B_i(t)$. Time-average occupancy $L_{q,i} = \frac{1}{T}\sum_{t=1}^T Q_i(t)$, mean queue wait $W_{q,i} = \frac{1}{N_i}\sum_{k=1}^{N_i} D_{i,k}$, and utilization $\hat{\rho}_i = \frac{1}{Tc_i}\sum_{t=1}^T B_i(t)$ use fixed seed \texttt{20260922}. Multi-config sweeps (E2, E3) use \emph{Common Random Numbers (CRN)} so all compared configurations share identical arrival streams.

\begin{table}[!t]
\centering
\caption{Layer-1 Analytical \& Network Accounting Verification ($10^7$ Cycles, Seed \texttt{20260922})}
\label{tab:validation}
\resizebox{\columnwidth}{!}{%
\begin{tabular}{@{}llrrr@{}}
\toprule
\textbf{Validation Check \& Configuration} & \textbf{Metric / Relation} & \textbf{Analytical} & \textbf{Simulated} & \textbf{Rel.\ Error} \\
\midrule
\multicolumn{5}{@{}l}{\textit{V1: Single-Station Deterministic Service ($Geo/D/1$)}} \\
$S_1{=}2$, $\lambda{=}0.30$ ($\rho{=}0.60$) & $L_q = \lambda^2/(1{-}2\lambda)$ & 0.2250 & 0.2260 & 0.46\% \\
$S_1{=}2$, $\lambda{=}0.30$ ($\rho{=}0.60$) & $W_q = \rho(S{-}1)/(2(1{-}\rho))$ & 0.7500 & 0.7530 & 0.41\% \\
$S_1{=}2$, $\lambda{=}0.45$ ($\rho{=}0.90$) & $L_q = \lambda^2/(1{-}2\lambda)$ & 2.0250 & 2.0324 & 0.37\% \\
$S_1{=}2$, worst stable ($\lambda{=}0.35$) & $\max_{\lambda \le 0.45} |L_q - \hat{L}_q|/L_q$ & 0.4083 & 0.4053 & 0.75\% \\
$S_1{=}3$, $\lambda{=}0.30$ ($\rho{=}0.90$) & $L_q = \lambda W_q$ ($S$-general) & 2.7000 & 2.7118 & 0.44\% \\
$S_1{=}3$, $\lambda{=}0.30$ ($\rho{=}0.90$) & $W_q = \rho(S{-}1)/(2(1{-}\rho))$ & 9.0000 & 9.0416 & 0.46\% \\
$S_1{=}3$, ratio vs.\ $S{=}2$-only form & $\hat{L}_q / [\lambda^2(S{-}1)/(1{-}\lambda S)]$ & 1.5000 & 1.5066 & 0.44\% \\
\midrule
\multicolumn{5}{@{}l}{\textit{V2: Single-Station Geometric Service ($Geo/Geo/1$, $E[S]{=}2$) \& Factor-2 Law}} \\
$E[S]{=}2$, $\lambda{=}0.30$ ($\rho{=}0.60$) & $L_q = 2\lambda^2/(1{-}2\lambda)$ & 0.4500 & 0.4505 & 0.11\% \\
$E[S]{=}2$, $\lambda{=}0.45$ ($\rho{=}0.90$) & $L_q = 2\lambda^2/(1{-}2\lambda)$ & 4.0500 & 4.0775 & 0.68\% \\
$E[S]{=}2$, worst stable ($\lambda{=}0.10$) & $\max_{\lambda \le 0.45} |L_q - \hat{L}_q|/L_q$ & 0.0250 & 0.0248 & 0.94\% \\
Factor-2 ratio ($\lambda{=}0.45$) & $\hat{L}_q^{\text{Geo}} / \hat{L}_q^{\text{D}}$ & 2.0000 & 2.0062 & 0.31\% \\
Factor-2 sweep ($\lambda \in [0.05, 0.45]$) & Range of $\hat{L}_q^{\text{Geo}} / \hat{L}_q^{\text{D}}$ & 2.0000 & $1.984\text{--}2.019$ & $\le 0.95\%$ \\
\midrule
\multicolumn{5}{@{}l}{\textit{V3: Full 3-Station Pipeline Accounting ($B{=}1, \tau{=}0, c{=}1$, $\lambda{=}0.30$)}} \\
Flow conservation ($Q_1, Q_2, Q_3$) & $\text{Departures} / \text{Arrivals}$ & 1.0000 & 0.9999997 & $3.3{\times}10^{-7}$ \\
Station utilizations ($\hat{\rho}_1, \hat{\rho}_3$) & $\rho_i = \lambda_i E[S_i]$ & 0.6000 & 0.5997 / 0.5999 & $\le 0.06\%$ \\
Per-node $W = L_q / \lambda$ relation & $\hat{W}_{q,i} / (\hat{L}_{q,i}/\hat{\lambda}_i)$ & 1.0000 & 1.0000000 & $1.5{\times}10^{-16}$ \\
Network total occupancy & $\hat{L}_{\text{tot}} / (\hat{\lambda} \hat{W}_{\text{e2e,qwait}})$ & 0.5591 & 0.5590868 & $3.3{\times}10^{-7}$ \\
\bottomrule
\end{tabular}%
}
\vspace{-2pt}
\end{table}

\begin{figure}[!t]
\centering
\begin{minipage}[t]{0.48\linewidth}
\centering
\incplot{../Project_Simulator/figures/fig_e1_validation_geo_d_1.pdf}{fig_e1_validation_geo_d_1.pdf}{\linewidth}
\end{minipage}\hfill
\begin{minipage}[t]{0.48\linewidth}
\centering
\incplot{../Project_Simulator/figures/fig_v2_factor2.pdf}{fig_v2_factor2.pdf}{\linewidth}
\end{minipage}
\vspace{-2pt}
\caption{Layer-1 verification. \emph{Left (V1):} $Geo/D/1$ ($S_1{=}2$) queue occupancy across horizons $T \in \{10^3, 10^5, 10^7\}$ cycles, converging onto $L_q = \lambda^2/(1{-}2\lambda)$ ($0.46\%$ error at $\lambda{=}0.30, 10^7$ cycles). \emph{Right (V2):} $Geo/Geo/1$ vs.\ $Geo/D/1$ at $10^7$ cycles, confirming the exact $2\times$ Pollaczek-Khinchine variance penalty.}
\label{fig:layer1_val}
\vspace{-3pt}
\end{figure}

\section{Layer-1 Analytical Verification}
We verify single-station and pipeline behavior across $\lambda \in \{0.05, \dots, 0.45, 0.49\}$ and $T \in \{10^3, 10^5, 10^7\}$ cycles. Additionally, \texttt{check\_conventions.py} confirms bit-exact integer counter equivalence with the course's \texttt{single\_queue\_simulation\_geo\_d\_1.m} across four $(S, \lambda, T)$ test cases.

\textbf{Deterministic Service ($Geo/D/1$) \& $S$-General Formula.}
From the waiting-time PGF $G_W(z) = (1 - \lambda S)/[1 - \lambda(1 + z + \dots + z^{S-1})]$, the exact $S$-general closed-form pair (\texttt{top\_script.m}) is:
\begin{equation}
W_q = \frac{\rho(S - 1)}{2(1 - \rho)}, \qquad L_q = \lambda W_q = \frac{\lambda^2 S(S - 1)}{2(1 - \lambda S)}.
\label{eq:geod1_general}
\end{equation}
During verification we proved that the Assignment-1 formula $L_q^{\text{A1}} = \lambda^2(S - 1)/(1 - \lambda S)$ matches our post-dispatch convention \emph{only} at $S{=}2$ ($S/2{=}1 \implies L_q = \lambda^2/(1 - 2\lambda)$), whereas at $S{=}3$ the exact occupancy is $S/2 = 1.5\times$ larger. Table~\ref{tab:validation} and Fig.~\ref{fig:layer1_val} (left) show that at $10^7$ cycles, our simulator matches Eq.~(\ref{eq:geod1_general}) within $0.75\%$ across $\lambda \le 0.45$ for $S_1{=}2$ and within $0.46\%$ for $S_1{=}3$ ($\lambda \le 0.30$, confirming the $1.5066\times$ ratio). At $\lambda{=}0.30$, $L_q$ relative error converges from $7.56\%$ ($10^3$ cycles) to $1.83\%$ ($10^5$) and $0.46\%$ ($10^7$).

\textbf{Geometric Service ($Geo/Geo/1$) \& V3 Accounting.}
For geometric service with $E[S]{=}2$ ($\text{Var}(S){=}2$), discrete-time Pollaczek-Khinchine gives $L_q^{\text{Geo}} = 2\lambda^2/(1 - 2\lambda) = 2L_q^{\text{D}}$. Table~\ref{tab:validation} and Fig.~\ref{fig:layer1_val} (right) confirm agreement within $0.94\%$ across $\lambda \le 0.45$ and an empirical ratio $\hat{L}_q^{\text{Geo}}/\hat{L}_q^{\text{D}} \in [1.9835, 2.0190]$ ($0.31\%$ error at $\lambda{=}0.45$). All 87 stable checks in \texttt{validation\_table.csv} pass with zero failures.

\begin{table}[!t]
\centering
\caption{Layer-2 Empirical Tesla T4 GPU Validation, Uncapped Predictor Comparison \& Multi-Model Batch-Scaling Sweeps}
\label{tab:gpu_empirical}
\resizebox{\columnwidth}{!}{%
\begin{tabular}{@{}lrrrrrr@{}}
\toprule
\multicolumn{7}{@{}l}{\textbf{Part A: Split Calibration ($B \in \{1,4,16\}$) \& Held-Out Validation ($B \in \{2,8,32\}$), 96-Token Cap}} \\
\textbf{Batch} & \textbf{Split} & \textbf{Colab} & \textbf{Kaggle} & \textbf{Power-Law} & \textbf{Latency} & \textbf{Measured / Pred.} \\
\textbf{Size $B$} & \textbf{Role} & \textbf{Mean (s)} & \textbf{Mean (s)} & \textbf{Pred.\ (s)} & \textbf{Error} & \textbf{Req/s (Error)} \\
\midrule
1 & Calib. & $3.868 \pm 0.29$ & 4.069 & 4.030 & -- & 0.259 / 0.248 (--) \\
4 & Calib. & $4.382 \pm 2.21$ & 3.826 & 4.036 & -- & 0.913 / 0.991 (--) \\
16 & Calib. & $3.879 \pm 0.29$ & 3.980 & 4.042 & -- & 4.125 / 3.959 (--) \\
\textbf{2} & \textbf{Held-Out} & $\mathbf{3.795 \pm 0.28}$ & \textbf{3.869} & \textbf{4.033} & \textbf{+6.26\%} & \textbf{0.527 / 0.496 (-5.89\%)} \\
\textbf{8} & \textbf{Held-Out} & $\mathbf{3.800 \pm 0.28}$ & \textbf{3.865} & \textbf{4.039} & \textbf{+6.29\%} & \textbf{2.105 / 1.981 (-5.92\%)} \\
\textbf{32} & \textbf{Held-Out} & $\mathbf{4.015 \pm 0.30}$ & \textbf{4.088} & \textbf{4.045} & \textbf{+0.74\%} & \textbf{7.970 / 7.912 (-0.74\%)} \\
\multicolumn{5}{@{}l}{\textbf{Held-Out Mean Absolute Percentage Error (MAPE):}} & \textbf{4.43\%} & \textbf{Throughput: 4.18\%} \\
\midrule
\multicolumn{7}{@{}l}{\textbf{Part B: Uncapped 512-Token Structural Sweep (30 Trials, $L_j \in [87, 498]$ Tokens)}} \\
\multicolumn{2}{@{}l}{\textbf{Predictor Model}} & \multicolumn{2}{r}{\textbf{Feature Correlation $r$}} & \multicolumn{2}{r}{\textbf{Held-Out Latency MAPE}} & \textbf{Throughput MAPE} \\
\midrule
\multicolumn{2}{@{}l}{Linear-in-$B$: $\hat{T}(B) = \beta_0 + \beta_1 B$} & \multicolumn{2}{r}{$r(T, B) = 0.7384$} & \multicolumn{2}{r}{$15.25\%$} & $17.60\%$ \\
\multicolumn{2}{@{}l}{\textbf{Max-Length: $\hat{T}(L_{\max}) = \gamma_0 + \gamma_1 \max_j L_j$}} & \multicolumn{2}{r}{$\mathbf{r(T, L_{\max}) = 0.9493}$} & \multicolumn{2}{r}{$\mathbf{9.14\%}$} & $\mathbf{10.75\%}$ \\
\midrule
\multicolumn{7}{@{}l}{\textbf{Part C: Multi-Model Measured T4 Sweeps (Short-QA, $B \in \{1, \dots, 32\}$), Four Model Families}} \\
\multicolumn{2}{@{}l}{\textbf{Model (Trials per $B$)}} & \textbf{FP16 Wt.\ / Peak} & \textbf{KV Cache} & \textbf{Batch Time (s)} & \textbf{Req/s (Gain)} & \textbf{Scaling} \\
\multicolumn{2}{@{}l}{} & \textbf{GB ($B{=}32$)} & \textbf{per Token} & $\mathbf{B{=}1 \to 32}$ & $\mathbf{B{=}1 \to 32}$ & \textbf{Exponent $\alpha$} \\
\midrule
\multicolumn{7}{@{}l}{\textit{96-token cap (controlled hardware benchmark)}} \\
\multicolumn{2}{@{}l}{\texttt{Qwen2.5-1.5B} (10)} & 3.0 / -- & 28 KB & $3.87 \to 4.02$ & $0.26 \to 7.97$ ($31\times$) & 0.001 \\
\multicolumn{2}{@{}l}{\texttt{Gemma-2-2B-it} (6)} & 5.2 / 5.7 & 104 KB & $4.79 \to 5.07$ & $0.21 \to 6.31$ ($30\times$) & 0.012 \\
\multicolumn{2}{@{}l}{\texttt{Llama-3.2-1B-Inst.} (10)} & 2.5 / 2.7 & 32 KB & $2.63 \to 2.87$ & $0.38 \to 11.14$ ($29\times$) & 0.022 \\
\multicolumn{7}{@{}l}{\textit{512-token ceiling (uncapped generation lengths)}} \\
\multicolumn{2}{@{}l}{\texttt{Llama-3.2-1B-Inst.} (5)} & 2.5 / 3.2 & 32 KB & $9.84 \to 22.15$ & $0.10 \to 1.45$ ($14\times$) & 0.20 \\
\multicolumn{2}{@{}l}{\texttt{Phi-3-mini-4k} (5)} & 7.6 / 14.5 & 384 KB & $5.58 \to 61.29$ & $0.18 \to 0.52$ ($2.9\times$) & 0.67 \\
\bottomrule
\end{tabular}%
}
\vspace{-2pt}
\end{table}

\section{Layer-2 Empirical GPU Analysis and Simulation Results}
To address Phase-1 Feedback~\#2 (\emph{stronger experiments, diverse models, and heterogeneous queries}), we executed hardware profiling on Nvidia Tesla T4 GPUs alongside multi-parameter simulator campaigns E2 and E3.

\subsection{Empirical GPU Grounding \& Split Calibration (Tesla T4, Qwen 2.5-1.5B)}
We profiled \texttt{Qwen/Qwen2.5-1.5B-Instruct} in FP16 (\texttt{sdpa} attention, KV-cache enabled) on Nvidia Tesla T4 GPUs ($15.36\text{ GB}$ VRAM, CUDA 12.8, PyTorch 2.11) across Google Colab and Kaggle. We swept static batch sizes $B \in \{1, 2, 4, 8, 16, 32\}$ with 10 trials per size (60 batch runs, 630 generated sequences) using CUDA-synchronized timers, plus 20 CPU tokenizer trials ($0.206\text{ ms}$ mean, $\text{SD}{=}0.283\text{ ms}$, range $[0.068, 1.369]\text{ ms}$).

\textbf{Out-of-Sample Power-Law Calibration.}
We split the batch sweep into a \emph{calibration set} $B_{\text{cal}} \in \{1, 4, 16\}$ and an unseen \emph{held-out test set} $B_{\text{test}} \in \{2, 8, 32\}$. Log-linear OLS regression on $B_{\text{cal}}$ yields:
\begin{equation}
T_{\text{batch}}(B) = a \cdot B^\alpha, \quad a = 4.0301\text{ s}, \quad \alpha = 0.001044.
\label{eq:power_law}
\end{equation}
As shown in Table~\ref{tab:gpu_empirical} (Part~A) and Fig.~\ref{fig:gpu_plots} (left), predicting on held-out sizes $B \in \{2, 8, 32\}$ achieves a \textbf{Batch Latency MAPE of $4.43\%$} (individual errors $+6.26\%, +6.29\%, +0.74\%$) and a \textbf{Throughput MAPE of $4.18\%$}. Because $\alpha \approx 0.001$ under the 96-token cap ($90\text{--}98\%$ of sequences hit 96 tokens), batch latency stays nearly flat ($\approx 3.80\text{--}4.01\text{ s}$) from $B{=}1$ to $32$, scaling throughput $30.8\times$ from $24.59\text{ tok/s}$ ($0.259\text{ req/s}$ at $B{=}1$) to $763.28\text{ tok/s}$ ($7.970\text{ req/s}$ at $B{=}32$). Cross-platform Kaggle vs.\ Colab means agree within $1.70\text{--}4.94\%$ across $B \in \{1, 2, 8, 16, 32\}$.

\begin{figure}[!t]
\centering
\begin{minipage}[t]{0.52\linewidth}
\centering
\incplot{../GPU_Measurements/colab_kaggle_sweep/sweep_20260920T132814Z_96c7e45f/colab_kaggle_comparison.png}{colab_kaggle_comparison.png}{\linewidth}
\end{minipage}\hfill
\begin{minipage}[t]{0.46\linewidth}
\centering
\incplot{../GPU_Measurements/NAME3_GPU_Analysis/geometric_cdf_comparison.png}{geometric_cdf_comparison.png}{\linewidth}
\end{minipage}
\vspace{-2pt}
\caption{Layer-2 Tesla T4 GPU measurements (\texttt{Qwen2.5-1.5B-Instruct}). \emph{Left:} Batch execution time and throughput across $B \in \{1,\dots,32\}$ on Colab and Kaggle with power-law calibration on $\{1,4,16\}$ and held-out validation on $\{2,8,32\}$ ($4.43\%$ MAPE). \emph{Right:} Uncapped 100-prompt empirical sequence-length CDF vs.\ fitted geometric CDF ($KS = 0.3965$).}
\label{fig:gpu_plots}
\vspace{-3pt}
\end{figure}

\subsection{Query Diversity, Straggler Length \& Non-Geometric Distributions}
To move beyond fixed 96-token caps and evaluate natural query diversity, we conducted two uncapped diagnostic studies with a 512-token ceiling:
\begin{enumerate}\setlength{\itemsep}{0pt}\setlength{\topsep}{1pt}
  \item \emph{100-Prompt Uncapped Length Diagnostic:} Across 100 diverse prompts, $0/100$ hit the 512-token cap. Generated lengths spanned $[83, 487]$ tokens (median $254.5$, mean $255.23$, $\text{SD}{=}74.57$, P90 $330.4$, P99 $442.5$ tokens; wall time $3.36\text{--}19.80\text{ s}$, mean $10.30\text{ s}$). Fig.~\ref{fig:gpu_plots} (right) shows that comparing the empirical CDF against a fitted memoryless geometric distribution yields $\mathbf{KS = 0.3965}$, proving that real LLM outputs exhibit a positive minimum-length floor and heavy concentration around $180\text{--}330$ tokens rather than memoryless step-1 decay.
  \item \emph{Uncapped Straggler Coupling ($r{=}0.9493$):} In static batching, a worker holds the batch until the \emph{longest sequence} $L_{\max} = \max_{j \le B} L_j$ completes. Across 30 uncapped batch trials ($B \in \{1,\dots,32\}$), batch time correlates moderately with $B$ ($r{=}0.7384$) only because $E[\max_{j \le B} L_j]$ grows from $232.8$ tokens ($B{=}1$) to $454.2$ tokens ($B{=}32$), whereas correlation with $L_{\max}$ is $\mathbf{r = 0.9493}$ ($\approx 40.3\text{ ms/token}$). As shown in Table~\ref{tab:gpu_empirical} (Part~B), conditioning on $L_{\max}$ cuts held-out latency MAPE on $B \in \{2, 8, 32\}$ from $\mathbf{15.25\%}$ (linear-in-$B$) to $\mathbf{9.14\%}$ and throughput MAPE from $\mathbf{17.60\%}$ to $\mathbf{10.75\%}$.
\end{enumerate}

\subsection{Multi-Model Batch-Scaling Sweeps \& Cycle Mapping}
Beyond the \texttt{Qwen2.5-1.5B-Instruct} calibration ($\approx 3.0\text{ GB}$ FP16 weights, HBM weight-read floor $\approx 9.4\text{ ms/step}$ at $320\text{ GB/s}$ vs.\ measured $\approx 40.3\text{ ms/step}$), Table~\ref{tab:gpu_empirical} (Part~C) reports measured T4 sweeps on three further model families (\texttt{Gemma-2-2B-it}, \texttt{Llama-3.2-1B-Instruct}, \texttt{Phi-3-mini-4k-instruct}) over the same $B \in \{1, \dots, 32\}$ with 5--10 trials per size. Under the 96-token cap, all three families reproduce the flat service law seen for Qwen: batch time grows by only $4\text{--}9\%$ from $B{=}1$ to $B{=}32$ ($\alpha{=}0.012$ for Gemma, $0.022$ for Llama), so request throughput rises by $\approx 29\text{--}31\times$, and measured per-step decode time ($\approx 27\text{ ms}$ Llama, $\approx 50\text{ ms}$ Gemma) sits $3.1\text{--}4.3\times$ above each model's weight-read floor ($7.7$ and $16.3\text{ ms}$). Without the cap, scaling is model-dependent: \texttt{Llama-3.2-1B} remains strongly sublinear ($\alpha{=}0.20$, $14\times$ throughput gain), whereas \texttt{Phi-3-mini} ($7.6\text{ GB}$ weights, $384\text{ KB}$ KV-cache/token) is close to linear ($\alpha{=}0.67$, only $2.9\times$ gain) and peaks at $14.5\text{ GB}$ of the $15.4\text{ GB}$ T4 at $B{=}32$ (two larger-batch runs on other workloads ran out of memory). The capped and uncapped rows are separate regimes and are not compared across groups: uncapped Llama batches are dominated by sequences reaching the 512-token ceiling. Finalizing the physical \texttt{cycle\_ms} mapping in \texttt{calibration\_parameters.json} is scheduled for Phase~3.

\begin{table}[!t]
\centering
\caption{Simulation Experiments E2 (Worker Scaling) \& E3 (Batching Trade-Offs, $\lambda{=}0.30, \delta{=}0.5, 10^7$ Cycles)}
\label{tab:e2_e3}
\resizebox{\columnwidth}{!}{%
\begin{tabular}{@{}rrrrrrrr@{}}
\toprule
\multicolumn{8}{@{}l}{\textbf{Part A: Experiment E2: Service Variability \& Worker Parallelism ($B{=}1, \tau{=}0, 10^6$ Cycles)}} \\
$\lambda$ & $c1\text{-D}$ $W_{\text{soj}}$ & $c1\text{-Geo}$ $W_{\text{soj}}$ & \textbf{Var.\ Premium} & $c2\text{-Geo}$ $W_{\text{soj}}$ & $L_{q,1}$ (All) & $L_{q,3}$ ($c1\text{-D}$) & $L_{q,3}$ ($c1 / c2\text{-Geo}$) \\
\midrule
0.10 & 4.123 & 4.303 & $+4.4\%$ & 4.122 & 0.0123 & 0.0000 & 0.0183 / 0.0001 \\
0.20 & 4.332 & 4.835 & $+11.6\%$ & 4.340 & 0.0664 & 0.0000 & 0.1004 / 0.0014 \\
0.30 & 4.754 & 5.871 & $+23.5\%$ & 4.770 & 0.2264 & 0.0000 & 0.3349 / 0.0044 \\
0.40 & 6.018 & 9.067 & $+50.7\%$ & 6.050 & 0.8074 & 0.0000 & 1.2185 / 0.0111 \\
\textbf{0.45} & \textbf{8.519} & \textbf{15.350} & $\mathbf{+80.2\%}$ & \textbf{8.549} & \textbf{2.0346} & \textbf{0.0000} & \textbf{3.0761 / 0.0142} \\
0.49 & 28.823 & 69.151 & $+139.9\%$ & 28.862 & 12.1750 & 0.0000 & 19.7793 / 0.0185 \\
\midrule
\multicolumn{8}{@{}l}{\textbf{Part B: Experiment E3: Batching Scheduler Trade-Off Grid ($\lambda{=}0.30, \delta{=}0.5, c{=}1, 10^7$ Cycles)}} \\
$B$ & $\tau$ & \textbf{Mean Fill $\bar{b}$} & $W_{\text{soj}}$ (cyc) & $W_{\text{qwait}}$ (cyc) & $L_{q,2}$ (req) & \textbf{$Q_3$ Util.\ $\hat{\rho}_3$} & \textbf{Throughput} \\
\midrule
\textbf{1} & \textbf{2 or 8} & \textbf{1.000} & \textbf{4.752} & \textbf{0.752} & \textbf{0.000} & \textbf{0.5999} & \textbf{0.3000} \\
2 & 2 & 1.574 & 6.752 & 2.023 & 0.381 & 0.4906 & 0.3000 \\
4, 8, 16 & 2 & 1.574 & 6.752 & 2.023 & 0.381 & 0.4906 & 0.3000 \\
2 & 8 & 1.953 & 7.360 & 2.384 & 0.489 & 0.4536 & 0.3000 \\
4 & 8 & 3.172 & 10.438 & 4.875 & 1.236 & 0.3224 & 0.3000 \\
\textbf{8} & \textbf{8} & \textbf{3.329} & \textbf{10.969} & \textbf{5.391} & \textbf{1.391} & \textbf{0.3063} & \textbf{0.3000} \\
\textbf{16} & \textbf{8} & \textbf{3.329} & \textbf{10.969} & \textbf{5.391} & \textbf{1.391} & \textbf{0.3063} & \textbf{0.3000} \\
\bottomrule
\end{tabular}%
}
\vspace{-2pt}
\end{table}

\begin{figure}[!t]
\centering
\incplot{../Project_Simulator/figures/fig_e2_e2e_delay.pdf}{fig_e2_e2e_delay.pdf}{0.68\linewidth}
\vspace{-2pt}
\caption{Experiment E2 ($10^6$-cycle CRN sweep, $B{=}1$): End-to-end sojourn delay ($W_{\text{soj}}$) vs.\ injection rate $\lambda$. Geometric decode variance at $c{=}1$ inflates delay by $80.2\%$ at $\lambda{=}0.45$ ($15.35$ vs.\ $8.52$ cycles), whereas $c{=}2$ parallel workers absorb the penalty ($8.55$ cycles).}
\label{fig:e2_delay}
\vspace{-3pt}
\end{figure}

\begin{figure}[!t]
\centering
\incplot{../Project_Simulator/figures/fig_e3_batching.pdf}{fig_e3_batching.pdf}{0.88\linewidth}
\vspace{-2pt}
\caption{Experiment E3 ($\lambda{=}0.30, \delta{=}0.5, 10^7$ cycles, CRN-paired): End-to-end sojourn $W_{\text{soj}}$, throughput, and mean batch fill $\bar{b}$ across $B \in \{1,2,4,8,16\}$ and $\tau \in \{2,8\}$. At $\tau{=}8$, fill and delay saturate between $B{=}4$ ($\bar{b}{=}3.17$) and $B{=}8$ ($\bar{b}{=}3.33$).}
\label{fig:e3_batching}
\vspace{-3pt}
\end{figure}

\subsection{Experiment E2: Service Variability \& Worker Parallelism}
Experiment E2 (\texttt{experiment\_e2\_pipeline.csv}) sets $Q_2$ to pass-through ($B{=}1, \tau{=}0$) and compares deterministic ($S_3{=}2$) vs.\ geometric ($E[S_3]{=}2$) GPU service across $c \in \{1, 2\}$ under CRN-paired arrivals. Table~\ref{tab:e2_e3} (Part~A) and Fig.~\ref{fig:e2_delay} show:
\textbf{(1) Structural Zero-Queue Property:} In all 22 deterministic runs ($c \in \{1,2\}$, all $\lambda$), $L_{q,3} = 0.0000$ identically: $Q_1$ ($S_1{=}2$) spaces departures by $\ge 2$ cycles, so a 2-cycle $Q_3$ service finishes before the next arrival dispatches.
\textbf{(2) Variability Premium at $c{=}1$:} Geometric $Q_3$ service builds a large queue ($L_{q,3}{=}3.076$ at $\lambda{=}0.45$), inflating sojourn by $+23.5\%$ at $\lambda{=}0.30$ ($4.754 \to 5.871$ cycles), $+80.2\%$ at $\lambda{=}0.45$ ($8.519 \to 15.350$ cycles), and $+139.9\%$ at $\lambda{=}0.49$ ($28.823 \to 69.151$ cycles).
\textbf{(3) Parallel Absorption ($c{=}2$):} Adding a second worker ($c2\text{-Geo}$) cuts $L_{q,3}$ at $\lambda{=}0.45$ to $0.0142$, reducing sojourn from $15.350$ to $8.549$ cycles (within $0.030$ cycles of deterministic $8.519$).

\subsection{Experiment E3: Static Batching Scheduler Trade-Offs}
Experiment E3 (\texttt{experiment\_e3\_batching.csv}) sweeps $B \in \{1, 2, 4, 8, 16\} \times \tau \in \{2, 8\}$ at $\lambda{=}0.30$ with sublinear batch service $S_3(b) = \lceil 2 + 0.5(b - 1) \rceil$ ($\delta{=}0.5, c{=}1$) over $10^6$ and $10^7$ cycles (agreeing within $\le 0.0038$ cycles / $0.29\%$). Table~\ref{tab:e2_e3} (Part~B) and Fig.~\ref{fig:e3_batching} demonstrate:
\textbf{(1) Delay vs.\ Utilization Pareto Trade-Off:} Unbatched serving ($B{=}1$) minimizes end-to-end sojourn ($W_{\text{soj}}{=}4.752$ cycles, $W_q{=}0.752$ cycles, $L_{q,2}{=}0$) but requires high GPU worker utilization ($\hat{\rho}_3{=}0.5999$). Batching at $\tau{=}8, B{\ge}8$ achieves mean fill $\bar{b}{=}3.329$ ($3.400$ cycles/batch), slashing GPU utilization by $\mathbf{48.9\%}$ ($\hat{\rho}_3: 0.5999 \to 0.3063$) while raising $Q_2$ wait ($W_{\text{soj}}{=}10.969$, $W_q{=}5.391$ cycles).
\textbf{(2) Short-Timer Invariance ($\tau{=}2$):} At $\tau{=}2$, the timer fires before $B$ requests accumulate behind $Q_1$, making $B \in \{2,4,8,16\}$ identical ($\bar{b}{=}1.574, W_{\text{soj}}{=}6.752, \hat{\rho}_3{=}0.4906$).
\textbf{(3) Saturation Knee ($\tau{=}8$):} Fill rises from $1.953$ ($B{=}2$) to $3.172$ ($B{=}4$) and saturates at $3.329$ for $B \in \{8, 16\}$ (identical rows), while throughput stays invariant at $0.3000\text{ req/cycle}$ across all 20 runs.

\section{Limitations, Conclusion, and Phase-3 Plan}
\textbf{Limitations \& Phase-3 Plan.} Our simulator models a feedforward chain with Bernoulli arrivals, static batching, and non-preemptive GPU service. In Phase~3, we will: (1)~extend sweeps in Table~\ref{tab:gpu_empirical} to summarization and code workloads, calibrating model-specific slopes $\delta$; (2)~finalize the physical \texttt{cycle\_ms} mapping; and (3)~incorporate network bounds (V4) and empirical $L_{\max}$ sampling at $Q_3$.

\textbf{Conclusion \& Code Availability.} This Phase-2 report resolves both Phase-1 feedback items via an explicit systems-to-queueing latency decomposition verified to $<0.95\%$ over $10^7$ cycles, multi-model Tesla T4 GPU sweeps across four model families (\texttt{Qwen2.5-1.5B} at $4.43\%$ held-out MAPE, \texttt{Gemma-2-2B}, \texttt{Llama-3.2-1B}, \texttt{Phi-3-mini-4k}), 100 uncapped diagnostics ($KS{=}0.3965$, $r{=}0.9493$, $9.14\%$ MAPE), and E2/E3 trade-off analyses. Full code and measurement traces are publicly available at \textcolor{blue}{\url{https://github.com/vaibhav3000/CSA-Modelling-Simulation-Project}} and \textcolor{blue}{\url{https://github.com/h2o-Noob/MS-Project}}.

\begin{thebibliography}{9}\setlength{\itemsep}{0.5pt}\setlength{\parskip}{0pt}\scriptsize
\bibitem{kwon2023vllm}
W.~Kwon \emph{et al.}, ``Efficient memory management for large language model serving with PagedAttention,'' in \emph{Proc. ACM SOSP}, 2023, pp. 611--626.
\bibitem{yu2022orca}
G.-I. Yu \emph{et al.}, ``Orca: A distributed serving system for transformer-based generative models,'' in \emph{Proc. USENIX OSDI}, 2022, pp. 521--538.
\bibitem{sheng2023flexgen}
Y.~Sheng \emph{et al.}, ``FlexGen: High-throughput generative inference of large language models with a single GPU,'' in \emph{Proc. ICML}, 2023, pp. 31094--31116.
\bibitem{agrawal2024vidur}
A.~Agrawal \emph{et al.}, ``Vidur: A large-scale simulation framework for LLM inference,'' in \emph{Proc. MLSys}, 2024.
\bibitem{mor2006discrete}
M.~Haviv, \emph{Queues: A Course in Queueing Theory}.\hskip 0.4em\relax Springer, 2013, vol. 191.
\bibitem{zheng2023response}
L.~Zheng \emph{et al.}, ``Judging LLM-as-a-judge with MT-Bench and Chatbot Arena,'' in \emph{Advances in NeurIPS}, vol. 36, 2023, pp. 46595--46623.
\end{thebibliography}

\end{document}
